\chapter{Tình Bạn Thật Và Giả}
\markboth{Tình Bạn Thật Và Giả}{Real and False Friendship}

\section{Bạn Đông Không Có Nghĩa Là Bớt Cô Đơn}
\begin{truyen}{Bạn Đông Không Có Nghĩa Là Bớt Cô Đơn}{Having Many Friends Does Not Mean Feeling Less Alone}
\chuhoa{M}{ột học sinh có thể ngồi giữa nhóm bạn rất đông mà vẫn thấy lạc lõng.}
\textit{(A student can sit among many friends and still feel out of place.)}

Bởi vì số lượng người quanh em không quyết định độ sâu của kết nối.
\textit{(Because the number of people around you does not determine the depth of connection.)}

Có những nhóm chỉ tồn tại nhờ nói chuyện phiếm, chụp ảnh chung, và giữ nhau khỏi bị lẻ loi.
\textit{(Some groups exist only for small talk, group photos, and avoiding loneliness.)}

Nhưng hễ ai gặp chuyện thật, cả nhóm biến mất.
\textit{(But when someone faces a real problem, the whole group disappears.)}

Một người bạn thật không cần xuất hiện suốt ngày.
\textit{(A real friend does not need to appear all the time.)}

Họ chỉ cần xuất hiện đúng lúc và không giả vờ.
\textit{(They only need to appear at the right time and not pretend.)}

\end{truyen}

\begin{baihoc}
Dạy học sinh đừng đo tình bạn bằng số lượng tương tác. Hãy đo bằng độ thật khi có khó khăn.
\textit{(Teach students not to measure friendship by the number of interactions. Measure it by honesty when trouble arrives.)}
\end{baihoc}

\begin{ghinhoanh}
\textbf{Short English to remember:}

Teach students not to measure friendship by the number of interactions.

Measure it by honesty when trouble arrives.

\end{ghinhoanh}

\begin{tuhocnhanh}
\textbf{Cau hoi tu hoc / Self-study questions}

1. Van de chinh la gi? \textit{What is the main problem?}

2. Cach xu ly tot hon la gi? \textit{What is a better action?}

3. Mot cau tieng Anh em muon nho la cau nao? \textit{Which English sentence do you want to remember?}
\end{tuhocnhanh}
\ngancach

\section{Bạn Tốt Không Luôn Chiều Em}
\begin{truyen}{Bạn Tốt Không Luôn Chiều Em}{A Good Friend Does Not Always Please You}
\chuhoa{C}{ó người gọi mình là bạn thân nhưng luôn cổ vũ mọi quyết định bốc đồng của mình.}
\textit{(Some people call themselves close friends and encourage every impulsive decision you make.)}

Nghe rất sướng.
\textit{(It feels great.)}

Nhưng người dám nói: ``Mày đang sai'' mới là người có nguy cơ mất lòng vì muốn giữ mày.
\textit{(But the one who dares say: ``You are wrong'' risks your displeasure because they want to keep you safe.)}

Tình bạn thật không chỉ là chia vui.
\textit{(Real friendship is not only about sharing fun.)}

Nó còn là can đảm kéo nhau lại khi một đứa đang trượt xuống.
\textit{(It is also the courage to pull each other back when one starts sliding down.)}

Học sinh cần hiểu sự khác nhau giữa người làm mình vui và người làm mình tốt hơn.
\textit{(Students need to understand the difference between someone who makes you feel good and someone who helps you become better.)}

Hai kiểu đó có khi không trùng nhau.
\textit{(Those two are not always the same person.)}

\end{truyen}

\begin{baihoc}
Một lời nhắc khó nghe đúng lúc đáng giá hơn mười câu chiều theo. Tình bạn tử tế không nuông chiều sự tự hủy.
\textit{(One hard truth at the right time is worth more than ten pleasing lies. Decent friendship does not indulge self-destruction.)}
\end{baihoc}

\begin{ghinhoanh}
\textbf{Short English to remember:}

One hard truth at the right time is worth more than ten pleasing lies.

Decent friendship does not indulge self-destruction.

\end{ghinhoanh}

\begin{tuhocnhanh}
\textbf{Cau hoi tu hoc / Self-study questions}

1. Van de chinh la gi? \textit{What is the main problem?}

2. Cach xu ly tot hon la gi? \textit{What is a better action?}

3. Mot cau tieng Anh em muon nho la cau nao? \textit{Which English sentence do you want to remember?}
\end{tuhocnhanh}
\ngancach

\section{Đừng Ở Lại Chỉ Vì Sợ Mất Nhóm}
\begin{truyen}{Đừng Ở Lại Chỉ Vì Sợ Mất Nhóm}{Do Not Stay Just Because You Fear Losing the Group}
\chuhoa{N}{hiều học sinh ở trong một nhóm bạn khiến mình mệt, chỉ vì sợ đi ra sẽ không còn ai.}
\textit{(Many students stay in draining friend groups only because they fear having no one if they leave.)}

Họ cười theo những câu nói ác, tham gia những trò khiến mình thấy bẩn, và gọi đó là hòa nhập.
\textit{(They laugh at cruel jokes, join in actions that make them feel dirty, and call it fitting in.)}

Nhưng cái giá của việc thuộc về một nhóm độc hại là đánh mất sự tôn trọng chính mình.
\textit{(But the price of belonging to a toxic group is the loss of self-respect.)}

Có những cô đơn sạch sẽ hơn rất nhiều so với một tình bạn làm em hỏng người.
\textit{(Some forms of loneliness are far cleaner than a friendship that corrupts you.)}

\end{truyen}

\begin{baihoc}
Hãy dạy học sinh rằng thuộc về không đáng giá nếu cái giá là phải phản bội điều đúng. Một nhóm không đáng để đổi lấy nhân cách.
\textit{(Teach students that belonging is not worth it if the price is betraying what is right. No group is worth your character.)}
\end{baihoc}

\begin{ghinhoanh}
\textbf{Short English to remember:}

Teach students that belonging is not worth it if the price is betraying what is right.

No group is worth your character.

\end{ghinhoanh}

\begin{tuhocnhanh}
\textbf{Cau hoi tu hoc / Self-study questions}

1. Van de chinh la gi? \textit{What is the main problem?}

2. Cach xu ly tot hon la gi? \textit{What is a better action?}

3. Mot cau tieng Anh em muon nho la cau nao? \textit{Which English sentence do you want to remember?}
\end{tuhocnhanh}
\ngancach

\section{Tình Bạn Cũng Cần Ranh Giới}
\begin{truyen}{Tình Bạn Cũng Cần Ranh Giới}{Friendship Also Needs Boundaries}
\chuhoa{C}{ó bạn thân mượn bài mãi, nhờ làm hộ mãi, trút hết cảm xúc tiêu cực lên mình mãi.}
\textit{(Some close friends keep borrowing homework, asking for favors, and dumping endless negative emotion on you.)}

Nhiều học sinh chịu đựng vì sợ nói ra sẽ mất bạn.
\textit{(Many students tolerate this because they fear speaking up will cost them the friendship.)}

Nhưng một mối quan hệ chỉ sống được bằng sự im lặng của một phía thì nó vốn đã lệch.
\textit{(But a relationship that survives only because one side stays silent is already distorted.)}

Nói ``tao không làm hộ mày nữa'' hay ``hôm nay tao không đủ sức nghe thêm'' không phải ích kỷ.
\textit{(Saying ``I will not do it for you anymore'' or ``I do not have the energy to hear more today'' is not selfish.)}

Đó là giữ đường biên để cả hai còn tôn trọng nhau.
\textit{(It is drawing a line so mutual respect can survive.)}

\end{truyen}

\begin{baihoc}
Tình bạn tốt không nuốt chửng cá nhân. Học sinh cần biết giúp người khác mà không biến mình thành chỗ bị khai thác.
\textit{(A good friendship does not swallow the individual. Students must learn to help others without becoming something to exploit.)}
\end{baihoc}

\begin{ghinhoanh}
\textbf{Short English to remember:}

A good friendship does not swallow the individual.

Students must learn to help others without becoming something to exploit.

\end{ghinhoanh}

\begin{tuhocnhanh}
\textbf{Cau hoi tu hoc / Self-study questions}

1. Van de chinh la gi? \textit{What is the main problem?}

2. Cach xu ly tot hon la gi? \textit{What is a better action?}

3. Mot cau tieng Anh em muon nho la cau nao? \textit{Which English sentence do you want to remember?}
\end{tuhocnhanh}
\ngancach

\section{Bạn Nhậu Không Phải Lúc Nào Cũng Là Bạn Đời}
\begin{truyen}{Bạn Nhậu Không Phải Lúc Nào Cũng Là Bạn Đời}{A Good-Time Companion Is Not Always a Lifelong Friend}
\chuhoa{C}{ó những người rất vui khi đi chơi cùng, nhưng biến mất ngay khi em gặp chuyện nặng.}
\textit{(Some people are very enjoyable to spend time with, but disappear the moment real trouble arrives.)}

Điều đó không làm họ thành người xấu hoàn toàn.
\textit{(That does not make them entirely bad people.)}

Chỉ là em cần gọi đúng vai trò của họ.
\textit{(It simply means you need to name their role correctly.)}

Không phải mọi người đi cùng một đoạn đều là người nên đặt ở tầng sâu nhất của lòng tin.
\textit{(Not everyone who walks with you for a while belongs in the deepest level of trust.)}

Học sinh hay đau vì nhầm lẫn giữa người hợp vui và người hợp đời.
\textit{(Students are often hurt because they confuse people who fit fun with people who fit life.)}

\end{truyen}

\begin{baihoc}
Tình bạn bền cần được phân loại tỉnh táo. Không phải ai vui vẻ với em cũng đủ sức giữ em khi có bão.
\textit{(Durable friendship requires sober classification. Not everyone cheerful with you is strong enough to hold you in a storm.)}
\end{baihoc}

\begin{ghinhoanh}
\textbf{Short English to remember:}

Durable friendship requires sober classification.

Not everyone cheerful with you is strong enough to hold you in a storm.

\end{ghinhoanh}

\begin{tuhocnhanh}
\textbf{Cau hoi tu hoc / Self-study questions}

1. Van de chinh la gi? \textit{What is the main problem?}

2. Cach xu ly tot hon la gi? \textit{What is a better action?}

3. Mot cau tieng Anh em muon nho la cau nao? \textit{Which English sentence do you want to remember?}
\end{tuhocnhanh}
\ngancach

\section{Người Bạn Giữ Miệng}
\begin{truyen}{Người Bạn Giữ Miệng}{The Friend Who Can Keep Their Mouth Closed}
\chuhoa{M}{ột học sinh kể bí mật cho bạn thân.}
\textit{(A student shared a secret with a close friend.)}

Hôm sau nửa lớp biết.
\textit{(The next day half the class knew.)}

Nhiều bạn trẻ tưởng thân thiết là đủ.
\textit{(Many young people think closeness is enough.)}

Nhưng một tình bạn không biết giữ miệng thì rất nguy hiểm.
\textit{(But a friendship that cannot hold confidence is dangerous.)}

Người đáng tin không chỉ là người tốt bụng.
\textit{(A trustworthy person is not only kind.)}

Họ còn là người biết chịu trách nhiệm với thông tin được giao.
\textit{(They are also someone who takes responsibility for what is entrusted to them.)}

Có những mối quan hệ tan không phải vì phản bội lớn, mà vì sự bất cẩn nhỏ lặp đi lặp lại.
\textit{(Some relationships collapse not through great betrayals, but through repeated small carelessness.)}

\end{truyen}

\begin{baihoc}
Muốn xây tình bạn thật, học sinh phải học cả phẩm chất kín miệng. Tin cậy là thứ được giữ bằng hành vi nhỏ.
\textit{(To build real friendship, students must learn the virtue of discretion. Trust is preserved through small acts.)}
\end{baihoc}

\begin{ghinhoanh}
\textbf{Short English to remember:}

To build real friendship, students must learn the virtue of discretion.

Trust is preserved through small acts.

\end{ghinhoanh}

\begin{tuhocnhanh}
\textbf{Cau hoi tu hoc / Self-study questions}

1. Van de chinh la gi? \textit{What is the main problem?}

2. Cach xu ly tot hon la gi? \textit{What is a better action?}

3. Mot cau tieng Anh em muon nho la cau nao? \textit{Which English sentence do you want to remember?}
\end{tuhocnhanh}
\ngancach

\section{Chơi Cùng Không Có Nghĩa Chung Giá Trị}
\begin{truyen}{Chơi Cùng Không Có Nghĩa Chung Giá Trị}{Spending Time Together Does Not Mean Sharing Values}
\chuhoa{H}{ai bạn rất hợp tính, cười nhiều, nói chuyện hợp gu.}
\textit{(Two friends got along easily, laughed often, and shared the same taste.)}

Nhưng đến lúc có việc liên quan đến đúng sai, họ lộ ra khác nhau hoàn toàn.
\textit{(But when a question of right and wrong arose, they proved completely different.)}

Một người xem quay cóp là bình thường.
\textit{(One considered cheating normal.)}

Người kia không làm được.
\textit{(The other could not accept it.)}

Từ đó họ bớt thân.
\textit{(After that they grew less close.)}

Không phải vì ghét nhau, mà vì tình bạn lâu dài cần nhiều hơn sự hợp vui.
\textit{(Not because they hated each other, but because long-term friendship requires more than easy enjoyment.)}

Giá trị sống khác biệt quá lớn sẽ sớm kéo hai người về hai phía.
\textit{(When values differ too greatly, life soon pulls two people in different directions.)}

\end{truyen}

\begin{baihoc}
Học sinh cần nhìn bạn không chỉ bằng cảm giác hợp, mà còn bằng nền giá trị. Cái sau mới quyết định đường dài.
\textit{(Students need to see friendship not only through chemistry, but through values. The latter decides the long road.)}
\end{baihoc}

\begin{ghinhoanh}
\textbf{Short English to remember:}

Students need to see friendship not only through chemistry, but through values.

The latter decides the long road.

\end{ghinhoanh}

\begin{tuhocnhanh}
\textbf{Cau hoi tu hoc / Self-study questions}

1. Van de chinh la gi? \textit{What is the main problem?}

2. Cach xu ly tot hon la gi? \textit{What is a better action?}

3. Mot cau tieng Anh em muon nho la cau nao? \textit{Which English sentence do you want to remember?}
\end{tuhocnhanh}
\ngancach

\section{Người Bạn Không Ghen Khi Em Sáng Lên}
\begin{truyen}{Người Bạn Không Ghen Khi Em Sáng Lên}{The Friend Who Does Not Resent Your Growth}
\chuhoa{C}{ó tình bạn ổn chừng nào cả hai còn ngang nhau.}
\textit{(Some friendships are stable only while both people remain equal.)}

Hễ một người tiến lên là bắt đầu có khó chịu ngầm.
\textit{(The moment one grows, quiet resentment begins.)}

Một người bạn tốt thật có thể chạnh lòng, nhưng cuối cùng vẫn chúc em tiến lên.
\textit{(A truly good friend may feel a sting, but in the end still wants your growth.)}

Người chỉ thích em khi em không vượt họ thực ra không yêu quý em.
\textit{(A person who likes you only when you do not surpass them does not really cherish you.)}

Họ chỉ thích vị trí của mình so với em.
\textit{(They only like their position relative to you.)}

Tình bạn trưởng thành phải chịu nổi sự thay đổi bậc thang của nhau.
\textit{(Mature friendship must be able to withstand shifts in one another’s position.)}

\end{truyen}

\begin{baihoc}
Hãy để ý ai vui thật khi em tiến bộ. Đó là dấu hiệu rất mạnh của một tình bạn lành.
\textit{(Notice who is genuinely glad when you improve. That is a strong sign of healthy friendship.)}
\end{baihoc}

\begin{ghinhoanh}
\textbf{Short English to remember:}

Notice who is genuinely glad when you improve.

That is a strong sign of healthy friendship.

\end{ghinhoanh}

\begin{tuhocnhanh}
\textbf{Cau hoi tu hoc / Self-study questions}

1. Van de chinh la gi? \textit{What is the main problem?}

2. Cach xu ly tot hon la gi? \textit{What is a better action?}

3. Mot cau tieng Anh em muon nho la cau nao? \textit{Which English sentence do you want to remember?}
\end{tuhocnhanh}
\ngancach

\section{Ở Bên Người Đang Xuống Dốc}
\begin{truyen}{Ở Bên Người Đang Xuống Dốc}{Staying Beside Someone Sliding Down}
\chuhoa{M}{ột học sinh bắt đầu nghỉ học nhiều, cáu gắt, điểm tụt.}
\textit{(A student started skipping school, becoming irritable, and falling in grades.)}

Nhiều bạn tránh dần vì thấy phiền.
\textit{(Many friends gradually avoided him because he felt burdensome.)}

Chỉ còn một đứa bạn vẫn nhắn: ``Tao không cứu được mày, nhưng tao không biến mất.''
\textit{(Only one friend kept messaging: ``I may not be able to save you, but I will not disappear.'')}

Câu đó không giải quyết ngay mọi thứ.
\textit{(That sentence did not solve everything immediately.)}

Nhưng nó giữ cho người đang xuống dốc biết rằng mình chưa bị bỏ hẳn.
\textit{(But it helped the struggling student know he had not been fully abandoned.)}

Đôi khi tình bạn thật không phải giải pháp.
\textit{(Sometimes real friendship is not the solution itself.)}

Nó là sợi dây để người ta chưa rơi quá sâu.
\textit{(It is the rope that keeps someone from falling too deep.)}

\end{truyen}

\begin{baihoc}
Không phải lúc nào bạn bè cũng sửa được vấn đề của nhau. Nhưng sự hiện diện tử tế đã là một giúp đỡ lớn.
\textit{(Friends cannot always fix each other’s problems. But decent presence is already a major help.)}
\end{baihoc}

\begin{ghinhoanh}
\textbf{Short English to remember:}

Friends cannot always fix each other’s problems.

But decent presence is already a major help.

\end{ghinhoanh}

\begin{tuhocnhanh}
\textbf{Cau hoi tu hoc / Self-study questions}

1. Van de chinh la gi? \textit{What is the main problem?}

2. Cach xu ly tot hon la gi? \textit{What is a better action?}

3. Mot cau tieng Anh em muon nho la cau nao? \textit{Which English sentence do you want to remember?}
\end{tuhocnhanh}
\ngancach

\section{Tình Bạn Không Phải Hợp Đồng Chịu Đựng}
\begin{truyen}{Tình Bạn Không Phải Hợp Đồng Chịu Đựng}{Friendship Is Not a Contract of Endless Endurance}
\chuhoa{C}{ó mối quan hệ kéo dài chỉ vì cả hai đều quen, chứ không còn tốt cho nhau.}
\textit{(Some relationships continue only because both people are used to them, not because they are still good for either person.)}

Một người luôn đòi hỏi, một người luôn nhịn.
\textit{(One always demands, the other always endures.)}

Rồi cả hai gọi đó là thân.
\textit{(Then both call it closeness.)}

Đến lúc người luôn nhịn kiệt sức và bước ra, người kia mới bảo là phản bội.
\textit{(When the one who always endured finally became exhausted and left, the other called it betrayal.)}

Không phải mọi kết thúc của tình bạn đều xấu.
\textit{(Not every ending of friendship is bad.)}

Có kết thúc là cách bảo toàn phần tử tế cuối cùng còn sót lại.
\textit{(Some endings preserve the last decent part still left.)}

\end{truyen}

\begin{baihoc}
Học sinh cần biết có những mối quan hệ nên rời đi đàng hoàng. Không phải giữ được hết mới là tốt.
\textit{(Students need to know that some relationships should be left with dignity. Keeping everything is not always goodness.)}
\end{baihoc}

\begin{ghinhoanh}
\textbf{Short English to remember:}

Students need to know that some relationships should be left with dignity.

Keeping everything is not always goodness.

\end{ghinhoanh}

\begin{tuhocnhanh}
\textbf{Cau hoi tu hoc / Self-study questions}

1. Van de chinh la gi? \textit{What is the main problem?}

2. Cach xu ly tot hon la gi? \textit{What is a better action?}

3. Mot cau tieng Anh em muon nho la cau nao? \textit{Which English sentence do you want to remember?}
\end{tuhocnhanh}
